\section{R\'esultats}

\begin{frame}
 \frametitle{Progression}
 \framesubtitle{}
 \begin{itemize}
   \item \textcolor{gray}{Introduction}
   \item \textcolor{gray}{G\'en\'eralit\'es sur l'interaction laser-atome
en champ fort}
   \item \textcolor{gray}{Nouvelle approche de r\'esolution 
dans l'espace des moments}
   \item \textbf{R\'esultats}
   \item Conclusion
 \end{itemize}
\end{frame}


\begin{frame}
 \frametitle{Validation du mod\`ele}
 \framesubtitle{Comparaison aux valeurs exactes}
\begin{textblock*}{400pt}(30pt,60pt)
\small \begin{itemize}
\item Energies propres
\end{itemize}
\begin{table}[p]
 \tiny\centering
\input{../tables/tab_eigval_l3_n100_1200_k03_beam.txt}
\end{table}
\end{textblock*}

\begin{textblock*}{400pt}(30pt,170pt)
\small\begin{itemize}
\item Fonctions propres
\end{itemize}
\begin{table}[p]
 \tiny\centering
 \input{../tables/tab_wavfun_nl_beam.txt}
\end{table}

\end{textblock*}



\begin{textblock*}{30pt}(250pt,35pt)
\tiny
\begin{equation*}
 \kappa_{\rm ex}=Z/n,\qquad \kappa=0.3,\qquad n_{\rm max}=1200, \qquad \ell_{\rm max}=3
\end{equation*}

\end{textblock*}


\begin{textblock*}{270pt}(150pt,60pt)
%\footnotesize
 \scriptsize
\begin{itemize}
 \item[$\bullet$] D\'ependance du nombre d'\'etats li\'es d\'ecrits aux param\`etres de la base
 \item[$\bullet$] Pr\'ecision de plus de 10 chiffres sur 60\% des \'etats-li\'es d\'ecrits
\end{itemize}

\end{textblock*}


\begin{textblock*}{270pt}(150pt,172pt)
%\footnotesize
 \scriptsize
\begin{itemize}
 \item[$\bullet$] Les pseudo-\'etats reproduisent les fonctions propres de l'hydrog\`ene
 \item[$\bullet$] Le recouvrement est sur 11 chiffres
\end{itemize}

\end{textblock*}

%\begin{itemize}
%  \item Comparaison des \'energies propres
%  \item Comparaison des vecteurs propres
%\end{itemize}
\end{frame}

\begin{frame}
 \frametitle{Validation du mod\`ele}
 \framesubtitle{Comparaison aux valeurs exactes - Forces d'oscillateur/Polarisabilit\'es}

\begin{textblock*}{400pt}(30pt,55pt)
\begin{table}[p]
 \tiny\centering
 \input{../tables/tab_oscstr_conf_mom_beam.txt}
\end{table}

\end{textblock*}


\begin{textblock*}{410pt}(140pt,185pt)
\begin{table}[p]
 \tiny\centering
 \input{../tables/tab_polarizab.txt}
\end{table}

\end{textblock*}


\begin{textblock*}{220pt}(20pt,185pt)
%\footnotesize
 \scriptsize
\begin{itemize}
 \item[$\bullet$] Excellent accord avec les r\'ef\'erences\footnote{
\Tiny
%\fontsize{4}{12} \selectfont
Bethe \textit{et al.}, Quantum Mechanics of One- and Two-Electron Systems.
Springer, Berlin, (1957) \\
\hspace{5.4mm}Harriman, Phys. Rev. 101, 594–598, (1956). \\
\hspace{5.4mm}Thu \textit{et al.}, J. Phys. B : At. Mol. Opt. Phys. 27, 4083–4094, (1994)}
 \item[$\bullet$] Validit\'e du mod\`ele pour diff\'erents hydrog\'eno\"ides 
 \item[$\bullet$] \textbf{Description tr\`es pr\'ecise de la structure}
\end{itemize}

\end{textblock*}

%\begin{itemize}
%  \item Table des forces d'oscillateur
%  \item R\`egle de sommation de Thomas-Reiche-Kuhn
%  \item Polarisabilit\'es 
%\end{itemize}


%\begin{textblock*}{300pt}(30pt,235pt)
%\scriptsize
%Bethe \textit{et al.}, Quantum Mechanics of One- and Two-Electron Systems.
%Springer, Berlin, (1957). \\
%Harriman, Phys. Rev. 101, 594–598, (1956).
%\end{textblock*}


\end{frame}


\begin{frame}
 \frametitle{Validation du mod\`ele}
 \framesubtitle{Validation en r\'egime haute fr\'equence - 1}

%\makebox[0pt][1]{
% \raisebox{-\totalheight}[10pt][10pt]{%
%   \includegraphics[width=75mm]{../figures/fig6.eps}}} %


\begin{textblock*}{270pt}(20pt,60pt)
\includegraphics[scale=.35]{../figures/fig6.eps}
\end{textblock*}

\begin{textblock*}{270pt}(230pt,55pt)
\includegraphics[scale=.4, trim=0 0 100 0]{../figures/fig5.eps}
\end{textblock*}


\begin{textblock*}{270pt}(30pt,205pt)
\begin{itemize}
\scriptsize
 \item Premiers tests \`a haute fr\'equence et 
en champ intense
 \item Excellent recouvrement de notre mod\`ele avec la MCDT\footnote{
\Tiny
Hamido \textit{et al.}, Phys. Rev. A 84, 013422, (2011).}
\end{itemize}
\end{textblock*}

%\begin{textblock*}{270pt}(285pt,225pt)
%\scriptsize
%Hamido \textit{et al.}, Phys. Rev. A 84, 013422, (2011).
%\end{textblock*}
%\begin{tikzpicture}
% %\node (0) {\includegraphics[width=55mm]{../figures/fig6.eps}};
% \node at (0,.5) {\includegraphics[width=75mm]{../figures/fig6.eps}};
% %\node [right=of 0] {\includegraphics[width=55mm]{../figures/fig6.eps}};
% \node at (7,-0){\includegraphics[width=70mm, trim=0 0 100 0]{../figures/fig5.eps}};
%\end{tikzpicture}

%\begin{picture}(0,0)
% \put(-0,-15){\includegraphics[width=40mm]{../figures/fig3.eps}}
%\end{picture}

%\begin{picture}(0,0)
% \put(-20,-80){\includegraphics[width=75mm]{../figures/fig6.eps}}
%\end{picture}

%\begin{picture}(0,0)
% \put(0,0){\includegraphics[width=90mm]{../figures/fig5.eps}}
%\end{picture}

% \begin{figure}
%  \begin{minipage}[p]{.4\textwidth}
%  \includegraphics[scale=.3]{../figures/fig2.eps} 
%  \end{minipage}
%  \begin{minipage}[p]{.4\textwidth}
%  \includegraphics[scale=.3]{../figures/fig2.eps} 
%  \end{minipage}
% \end{figure}

%\begin{itemize}
%  \item Courbes de Ongonwou \textit{et al.} 
%  
%\end{itemize}
\end{frame}



\begin{frame}
 \frametitle{Validation du mod\`ele}
 \framesubtitle{Validation en r\'egime haute fr\'equence - 2}


\begin{textblock*}{270pt}(30pt,60pt)
\includegraphics[scale=.30]{../figures/fig1.eps}
\end{textblock*}

\begin{textblock*}{270pt}(230pt,50pt)
\includegraphics[scale=.35, trim=0 0 100 0]{../figures/fig2.eps}
\end{textblock*}

%\begin{tikzpicture}
% %\node (0) {\includegraphics[width=55mm]{../figures/fig6.eps}};
% \node at (0,.5) {\includegraphics[width=65mm]{../figures/fig2.eps}};
% %\node [right=of 0] {\includegraphics[width=55mm]{../figures/fig6.eps}};
% \node at (7,-0){\includegraphics[width=60mm, trim=0 0 100 0]{../figures/fig1.eps}};
%\end{tikzpicture}

%\begin{itemize}
%  \item Courbes de Ongonwou \textit{et al.} (suite)
%\end{itemize}

\begin{textblock*}{270pt}(30pt,190pt)
\begin{itemize}
\scriptsize
 \item Accord qualitatif (positions et amplitudes des pics ATI)
 \item D\'esaccord progressif avec l'augmentation de la dur\'ee d'interaction 
 \item D\'esaccord quantitatif tr\`es net avec la MPS\footnote{
\Tiny
Tetchou Nganso \textit{et al.}, Phys. Rev. A 83, 013401, (2011)
}
\end{itemize}
\end{textblock*}

\end{frame}


\begin{frame}
 \frametitle{Analyse des r\'esultats en r\'egime basse fr\'equence}
 \framesubtitle{Pr\'esentation des diff\'erents signaux laser utilis\'es}

\begin{textblock*}{270pt}(30pt,60pt)
\includegraphics[scale=.23]{../figures/chap2/il1_il2_2c.eps}
\end{textblock*}

\begin{textblock*}{270pt}(180pt,60pt)
\includegraphics[scale=.23, trim=0 0 100 0]{../figures/chap2/il2_2c_0_ps42.eps}
\end{textblock*}

\begin{textblock*}{270pt}(320pt,60pt)
\includegraphics[scale=.23, trim=0 0 100 0]{../figures/chap2/il1_il3_8c.eps}
\end{textblock*}


\begin{textblock*}{500pt}(30pt,153pt)
\scriptsize
 \begin{itemize}
  \item IL1\footnote{\Tiny 
Tetchou Nganso, Ph.D. dissertation, Université de Douala, 2011 ;
Telnov \textit{et al.}, Phys. Rev. A 100, (4), (2019).}\hspace{1mm} $A(t)=A_0 \sin^2(\frac{\pi t}{\tau})\sin(\omega t + \phi)$
 \hspace{3mm} $n_c=2,\;\phi=0,\,\frac{\pi}{4},\,\frac{\pi}{2} ;
\quad n_c=8,\; I({10^{14}\,\rm W/cm^2}) = 1.3, 3, 5 ;
\quad n_c=20$%,\; {\rm H}(1s),  {\rm He^+}(2sp)$
  \item IL2\footnote{\Tiny 
Chen \textit{et al.}, Phys. Rev. A 86, 013410, (2012).}\hspace{2mm} $\mathcal{E}(t)=\mathcal{E}_0 \sin^2(\frac{\pi t}{\tau})\sin(\omega t + \phi)$
 \hspace{3mm} $n_c=2, 5, 10,
 \quad\phi=0, \frac{\pi}{4}, \frac{\pi}{2}$
  \item IL3\footnote{\Tiny 
Arb\'o \textit{et al.} Phys. Rev. Lett.96, (14), (2006).}\hspace{2mm} $\mathcal{E}(t)=\mathcal{E}_0 \sin^2(\frac{\pi t}{\tau})\cos(\omega t)$
\hspace{8mm} $n_c=8,\; I\left({10^{14}\,\rm W/cm^2}\right) = 
0.5, 1, 2 $
 \end{itemize}
\end{textblock*}



%\begin{textblock*}{500pt}(30pt,153pt)
%\tiny
% \begin{itemize}
%  \item IL1\footnote{\Tiny 
%Tetchou Nganso, Ph.D. dissertation, Université de Douala, 2011 ;
%Telnov \textit{et al.}, Phys. Rev. A 100, (4), (2019).}\hspace{1mm} $A(t)=A_0 \sin^2(\frac{\pi t}{\tau})\sin(\omega t + \phi)$
% \hspace{3mm} $n_c=2,\;\phi=\left\lbrace\begin{matrix}
%0 \\ 
%\pi/4 \\
%\pi/2
%\end{matrix}\right.
%\quad n_c=8,\; I({\rm W/cm^2}) = \left\lbrace\begin{matrix}
%1.3\\ 
%3\\ 
%5
%\end{matrix}\right.
%\quad n_c=20,\; Z = \left\lbrace\begin{matrix}
%1({\rm H})\\ 
%2({\rm He^+}) 
%\end{matrix}\right.$
%  \item IL2\footnote{\Tiny 
%Chen \textit{et al.}, Phys. Rev. A 86, 013410, (2012).}\hspace{2mm} $\mathcal{E}(t)=\mathcal{E}_0 \sin^2(\frac{\pi t}{\tau})\sin(\omega t + \phi)$
% \hspace{3mm} $n_c=\left\lbrace\begin{matrix}
%2 \\ 
%5 \\
%10
%\end{matrix}\right.;
% \qquad\phi=\left\lbrace\begin{matrix}
%0 \\ 
%\pi/4 \\
%\pi/2
%\end{matrix}\right.$
%  \item IL3\footnote{\Tiny 
%Arb\'o \textit{et al.} Phys. Rev. Lett.96, (14), (2006).}\hspace{2mm} $\mathcal{E}(t)=\mathcal{E}_0 \sin^2(\frac{\pi t}{\tau})\cos(\omega t)$
%\hspace{3mm} $n_c=8,\; I\left({10^{14}\,\rm W/cm^2}\right) = \left\lbrace\begin{matrix}
%0.5\\ 
%1\\ 
%2 
%\end{matrix}\right.$
% \end{itemize}
%\end{textblock*}

\begin{textblock*}{200pt}(280pt,200pt)
\scriptsize
\begin{itemize}
 \item[$\bullet$] R\'esultats relatifs \`a 22 runs
 \item[$\bullet$] Plus d'une semaine requise par run
\end{itemize}
\end{textblock*}

\end{frame}


\begin{frame}
 \frametitle{Analyse des r\'esultats en r\'egime basse fr\'equence}
 \framesubtitle{Populations des \'etats li\'es et probabilit\'e d'ionisation - 1}

%%\visible<2->{
%\begin{tikzpicture}
% \node at (0,6) {\includegraphics[width=.4\linewidth]{../figures/pop_time_states_ion_stacked_diego2006.eps}};
% \node at (5,6.4){\includegraphics[width=.30\linewidth, trim=0 0 100 0]{../figures/pop_time_states_ion_stacked_tetchou2010.eps}};
%\end{tikzpicture}
%%}

%\begin{itemize}
%  \item Evolution des populations dans le temps (courbes)
%\end{itemize}

\begin{textblock*}{270pt}(30pt,55pt)
\includegraphics[scale=.23]{../figures/chap5/pop_time_states_ion_stacked_diego2006.eps}
\end{textblock*}

\begin{textblock*}{270pt}(170pt,55pt)
\includegraphics[scale=.23, trim=0 0 100 0]{../figures/chap5/pop_time_states_ion_stacked_tetchou2010.eps}
\end{textblock*}


\begin{textblock*}{270pt}(30pt,220pt)
\scriptsize
\begin{itemize}
 \item Excitation/ionisation cons\'equente du syst\`eme 
seulement dans le r\'egime tunnel 
 \item Confirmation de l'ionisation quasi-totale \`a l'approche de 
l'intensit\'e de saturation
\end{itemize}
\end{textblock*}


\begin{textblock*}{270pt}(250pt,220pt)
\scriptsize
\begin{equation}
 I_s = 8 \times 10^{14}\,\rm W/cm^2
\end{equation}
\end{textblock*}

\end{frame}



\begin{frame}
 \frametitle{Analyse des r\'esultats en r\'egime basse fr\'equence}
 \framesubtitle{Populations des \'etats li\'es et probabilit\'e d'ionisation - 2}


\begin{textblock*}{270pt}(200pt,55pt)
\includegraphics[scale=.27]{../figures/chap5/pop_l_n_int1_cyc_cep_chen2012.eps}
\end{textblock*}

\begin{textblock*}{270pt}(-20pt,60pt)
\begin{table}
\tiny\centering
\scalebox{.8}{
\input{../tables/chen2012_tab_2_nopc2_ps2_i1.txt}}
\end{table}
\end{textblock*}


\begin{textblock*}{270pt}(-20pt,130pt)
\begin{table}
\tiny\centering
\scalebox{.8}{
\input{../tables/chen2012_tab_3_nopc10_ps2_i1.txt}}
\end{table}
\end{textblock*}


\begin{textblock*}{290pt}(20pt,205pt)
\scriptsize
\begin{itemize}
 \item Accord avec la m\'ethode spectrale de l'auteur\footnote{\Tiny 
Chen \textit{et al.}, Phys. Rev. A 86, 013410, (2012).}
pour les impulsions br\`eves
 \item Accord avec la m\'ethode des grilles pour les impulsions longues
\end{itemize}
\end{textblock*}

\end{frame}




\begin{frame}
 \frametitle{Analyse des r\'esultats en r\'egime basse fr\'equence}
 \framesubtitle{Populations des \'etats li\'es et probabilit\'e d'ionisation - 3}

%\begin{tikzpicture}
% \node at (0,6) {\includegraphics[width=.45\linewidth]{../figures/chap5/pop_3d_nl_phase_nc2_tetchou2010_chen2012.eps}};
% \node at (5,6) {\includegraphics[height=.3\linewidth,width=.4\linewidth]{../figures/chap5/pop_l_n_intens_diego2006_bar.eps}};
% \node at (9.5,6.3){\includegraphics[width=.3\linewidth, trim=0 0 0 0]{../figures/chap5/pop_3d_nl_telnov2019.eps}};
%\end{tikzpicture}
%
% \begin{itemize}
%   \item Populations finales $P_\ell$ des sous-\'etats 
%et $P_n$ \'etats (courbes)
%   \item Populations finales $P_{n\ell}$ des sous-\'etats
% \end{itemize}


\begin{textblock*}{270pt}(0pt,55pt)
\includegraphics[scale=.27]{../figures/chap5/pop_3d_nl_phase_nc2_tetchou2010_chen2012.eps}
\end{textblock*}

\begin{textblock*}{270pt}(134pt,65pt)
\includegraphics[scale=.2]{../figures/chap5/pop_l_n_intens_diego2006_bar.eps}
\end{textblock*}

\begin{textblock*}{270pt}(334pt,55pt)
\includegraphics[scale=.2]{../figures/chap5/pop_3d_nl_telnov2019.eps}
\end{textblock*}

\begin{textblock*}{290pt}(20pt,190pt)
\scriptsize
\begin{itemize}
 \item Forte d\'ependance des populations \`a la phase pour les impulsions br\`eves
 \item Forte d\'ependance des populations aux intensit\'es pour les impulsions longues
 \item D\'esaccord avec la loi $\ell=n-1$\footnote{\Tiny Piraux \textit{et al.} Phys. Rev.
A 96(4), (2017).} liant le niveau $n$ et le sous-niveau $\ell$ les plus peupl\'es 
 \item Similarit\'e du comportement de ${\rm He^+}(2p)$ \`a celui de ${\rm H}(1s)$ 
\end{itemize}
\end{textblock*}

\end{frame}



\begin{frame}
 \frametitle{Analyse des r\'esultats en r\'egime basse fr\'equence}
 \framesubtitle{Densit\'es d'\'energie - 1}

%\framezoom<1><2>[border](1cm,-3cm)(6cm,5cm)
%\begin{tikzpicture}
% \node at (0,6) {\includegraphics[width=.5\linewidth]{../figures/chap5/fig_time_pz_8c_e0_diego2006.jpg}};
%\end{tikzpicture}

\begin{textblock*}{270pt}(10pt,60pt)
\includegraphics[scale=.27]{../figures/chap5/fig_time_pz_8c_e0_diego2006.jpg}
\end{textblock*}


\begin{textblock*}{240pt}(220pt,90pt)
\footnotesize
 \begin{itemize}
   \item Visualisation directe de l'ionisation\footnote{\Tiny
A. de Bohan, Ph.D. dissertation, Universit\'e Catholique de Louvain, 2000.}
   \item Oscillation en opposition de phase de la densit\'e avec le champ 
$\boldsymbol{A}(t)$
   \item Ionisation d'autant plus rapide que l'intensit\'e est \'elev\'ee
   \item Pr\'esence de raies verticales fixes par rapport \`a l'axe $p_z$
donnant les positions des pics ATI
 \end{itemize}
\end{textblock*}


\begin{textblock*}{240pt}(55pt,220pt)
\scriptsize
 $\gamma \gg 1$ % = 1.3263$
\hspace{12mm}
 $\gamma \sim 1$ % = 0.9434$
\hspace{12mm}
 $\gamma \ll 1$ % = 0.6667$
\end{textblock*}

%\begin{itemize}
%  \item Evolution dans le temps des densit\'es d'\'energie
%\end{itemize}

\end{frame}


\begin{frame}
 \frametitle{Analyse des r\'esultats en r\'egime basse fr\'equence}
 \framesubtitle{Densit\'es d'\'energie - 2}

%\begin{tikzpicture}
% \node at (0,0) {\includegraphics[width=.7\linewidth]{../figures/chap5/energy_density_tetchou2010_chen2012.eps}};
% \node at (7,.5) {\includegraphics[width=.32\linewidth]{../figures/chap5/energy_density_phase_nc2510_0_chen2012.eps}};
%\end{tikzpicture}

\begin{textblock*}{270pt}(30pt,55pt)
\includegraphics[scale=.24]{../figures/chap5/energy_density_tetchou2010_chen2012.eps}
\end{textblock*}

\begin{textblock*}{270pt}(270pt,55pt)
\includegraphics[scale=.24]{../figures/chap5/energy_density_phase_nc2510_0_chen2012.eps}
\end{textblock*}

\begin{textblock*}{400pt}(20pt,220pt)
\scriptsize
 \begin{itemize}
   \item Asym\'etrie gauche/droite manifeste pour les impulsions br\`eves
   \item Disparition des effets de phases avec l'augmentation 
de la dur\'ee de l'interaction
 \end{itemize}
\end{textblock*}

% \begin{itemize}
%   \item Densit\'es d'\'energie avant et arri\`ere \`a la fin
%de l'interaction
% \end{itemize}
\end{frame}


\begin{frame}
 \frametitle{Analyse des r\'esultats en r\'egime basse fr\'equence}
 \framesubtitle{Spectres en \'energie - 1}
%% \putat{-0}{-120}{\includegraphics[scale=0.3]{../figures/chap5/energy_distrib_dif_pis2_nc2510_chen2012.eps}}
%
%\begin{tikzpicture}
% \node at (0,0) {\includegraphics[width=.5\linewidth]{../figures/chap5/energy_distrib_dif_pis2_nc2510_chen2012.eps}};
% \node at (7,.5) {\includegraphics[width=.45\linewidth]{../figures/chap5/energy_distrib_dif_intens_diego2006.eps}};
%\end{tikzpicture}
% 
% \begin{itemize}
%   \item Spectres en \'energie
%   \item Spectres en \'energie dans les angles solides 
%d'ouverture $\pi/6$ et $\pi/2$ vers l'avant et vers l'arri\`ere
% \end{itemize}

\begin{textblock*}{270pt}(40pt,55pt)
\includegraphics[scale=.26]{../figures/chap5/energy_distrib_dif_pis2_nc2510_chen2012.eps}
\end{textblock*}

\begin{textblock*}{270pt}(240pt,55pt)
\includegraphics[scale=.25]{../figures/chap5/energy_distrib_dif_intens_diego2006.eps}
\end{textblock*}

%\begin{textblock*}{90pt}(350pt,100pt)
%\scriptsize
%Position des pics
%\begin{equation*}
% E = (n+s)\omega - (I_p + U_p) \label{eqn:ati_cons_law}.
%\end{equation*}
%\end{textblock*}

\begin{textblock*}{350pt}(30pt,210pt)
\scriptsize
\begin{itemize}
 \item Augmentation g\'en\'erale de la probabilit\'e d'\'emission ATI 
avec la dur\'ee de l'impulsion
 \item Pas de variation de la probabilit\'e d'\'emission directe avec 
la dur\'ee de l'impulsion
 \item Augmentation "uniforme" de la probabilit\'e d'\'emission ATI 
avec l'intensit\'e de l'impulsion
\end{itemize}
\end{textblock*}


\end{frame}


\begin{frame}
 \frametitle{Analyse des r\'esultats en r\'egime basse fr\'equence}
 \framesubtitle{Spectres en \'energie - 2}
%% \putat{-0}{-120}{\includegraphics[scale=0.3]{../figures/chap5/energy_distrib_dif_pis2_nc2510_chen2012.eps}}
%
%\begin{tikzpicture}
% \node at (0,0) {\includegraphics[width=.6\linewidth]{../figures/chap5/fig_energy_distrib_h1s_he2s_i05_20c_iwf_telnov.eps}};
% \node at (7,.5) {\includegraphics[width=.50\linewidth]{../figures/chap5/fig_energy_distrib_he2s_he2p_i08_20c_iwf_telnov.eps}};
%\end{tikzpicture}
% 
% \begin{itemize}
%   \item Spectre en \'energie
%   \item Spectre en \'energie dans les angles solides 
%d'ouverture $\pi/6$ et $\pi/2$ vers l'avant et vers l'arri\`ere
% \end{itemize}

\begin{textblock*}{270pt}(30pt,55pt)
\includegraphics[scale=.35]{../figures/chap5/fig_energy_distrib_h1s_he2s_i05_20c_iwf_telnov.eps}
\end{textblock*}

\begin{textblock*}{270pt}(230pt,55pt)
\includegraphics[scale=.35]{../figures/chap5/fig_energy_distrib_he2s_he2p_i08_20c_iwf_telnov.eps}
\end{textblock*}

%\begin{textblock*}{110pt}(350pt,100pt)
%\scriptsize
%Position des pics
%\begin{equation*}
% E = (n+s)\omega - (I_p + U_p) \label{eqn:ati_cons_law}.
%\end{equation*}
%\end{textblock*}

\begin{textblock*}{390pt}(30pt,210pt)
\scriptsize
\begin{itemize}
 \item Diminution de l'amplitude des pics ATI avec leur ordre pour 
${\rm H}(1s)$ et ${\rm He^+}(2p)$
avec la dur\'ee de l'impulsion
 \item Pas d'asym\'etrie gauche/droite pour les impulsions longues ($n_c=20$)
 \item Augmentation non lin\'eaire de l'amplitude des pics avec l'intensit\'e 
pour ${\rm He^+}(2s)$
\end{itemize}
\end{textblock*}

\end{frame}


\begin{frame}
 \frametitle{Analyse des r\'esultats en r\'egime basse fr\'equence}
 \framesubtitle{Distributions en moment}

%\begin{tikzpicture}
% \node at (0,0) {\includegraphics[width=.6\linewidth]{../figures/chap5/mom_pnpz_distrib_tetchou2010_chen2012.eps}};
% \node at (7,0) {\includegraphics[width=.3\linewidth]{../figures/chap5/mom_pnpz_distrib_phase_nc2510_chen2012.eps}};
%% \node at (9,0) {\includegraphics[width=.3\linewidth]{../figures/chap5/mom_pnpz_distrib_intens_diego2006.eps}};
%\end{tikzpicture}
% 
% \begin{itemize}
%   \item Distributions en moment
%   \item Distributions angulaires
% \end{itemize}

\begin{textblock*}{270pt}(10pt,55pt)
\includegraphics[scale=.22]{../figures/chap5/mom_pnpz_distrib_tetchou2010_chen2012.eps}
\end{textblock*}

\begin{textblock*}{270pt}(210pt,55pt)
\includegraphics[scale=.22]{../figures/chap5/mom_pnpz_distrib_phase_nc2510_chen2012.eps}
\end{textblock*}

\begin{textblock*}{270pt}(330pt,55pt)
\includegraphics[scale=.23]{../figures/chap5/mom_pnpz_distrib_intens_diego2006.eps}
\end{textblock*}


%\begin{textblock*}{110pt}(350pt,100pt)
%\scriptsize
%Position des pics
%\begin{equation*}
% E = (n+s)\omega - (I_p + U_p) \label{eqn:ati_cons_law}.
%\end{equation*}
%\end{textblock*}

\begin{textblock*}{390pt}(30pt,210pt)
\scriptsize
\begin{itemize}
 \item Structure en \'eventail peu dense dans la r\'egion des basses \'energies
 \item Les stries marquent la position des pics ATI
 \item Disparition des effets de phases avec la dur\'ee de l'impulsion
%\item Augmentation de la probabilit\'e d'ionisation avec la dur\'ee  
\end{itemize}
\end{textblock*}

\end{frame}


\begin{frame}
 \frametitle{Analyse des r\'esultats en r\'egime basse fr\'equence}
 \framesubtitle{Distributions en impulsions transversale et longitudinale}


%\begin{tikzpicture}
% \node at (0,0) {\includegraphics[width=.33\linewidth]{../figures/chap5/momx_pz_distrib_end_phase_tetchou2010_chen2012.eps}};
% \node at (5,0) {\includegraphics[width=.35\linewidth]{../figures/chap5/momx_pz_distrib_end_intens_diego2006.eps}};
% \node at (10,0) {\includegraphics[width=.37\linewidth]{../figures/chap5/momx_pn_distrib_end_intens_diego2006.eps}};
%\end{tikzpicture}
%
% \begin{itemize}
%   \item Distributions en impulsion transversale
%   \item Distributions impulsion longitudinale
% \end{itemize}

\begin{textblock*}{270pt}(30pt,75pt)
\includegraphics[scale=.22]{../figures/chap5/momx_pz_distrib_end_phase_tetchou2010_chen2012.eps}
\end{textblock*}

\begin{textblock*}{270pt}(150pt,80pt)
\includegraphics[scale=.22]{../figures/chap5/momx_pz_distrib_end_intens_diego2006.eps}
\end{textblock*}

\begin{textblock*}{270pt}(280pt,80pt)
\includegraphics[scale=.20]{../figures/chap5/momx_pn_distrib_end_intens_diego2006.eps}
\end{textblock*}

\begin{textblock*}{300pt}(150pt,190pt)
\scriptsize
\begin{itemize}
 \item Les effets de phase sont port\'es dans la direction longitudinale
 \item Disparition des effets de phase dans la direction longitudinale avec 
l'intensit\'e
 \item Augmentation des \'energies des photo\'electrons ATI avec l'intensit\'e
\end{itemize}
\end{textblock*}


\end{frame}


